% GENERATED FILE. Edit canonical lessons, then run npm run generate:books.
% canonical-source-hash: fnv1a64:7a53c4d922b21f94
% canonical-lessons: BN-C389-chalok, BN-C389-kondaktor, BN-C389-helmet, BN-C389-ishara-kora, BN-C389-haka

\chapter{Drivers and Conductors}
\label{ch:bn-drivers-and-conductors}
\hlchaptermodality{389}

\begin{chapteropening}
\textbf{By the end of this chapter:} \emph{I can talk with a driver and a conductor, signal a stop and wear a helmet.}

Complete the last lesson of chapter 389: I can talk with a driver and a conductor, signal a stop and wear a helmet.
\end{chapteropening}

\section[\texorpdfstring{\bn{চালক} (chālôk)}{চালক (chālôk)}]{\bn{চালক} (chālôk) — a driver}

\label{lesson:BN-C389-chalok}



\begin{quote}
Before the new one: say the Bengali for a van, then the Bengali for to race.
\end{quote}

\subsection*{You'll want to know: \bn{চালক}}
\textbf{\bn{চালক}} — \emph{chālôk} — \textquotedblleft{}a driver\textquotedblright{}.

\begin{quote}
\bn{বাসের} \bn{চালক} \bn{হঠাৎ} \bn{বাস} \bn{থামালেন।}
\end{quote}

\begin{quote}
\emph{bāser chālôk hôṭhāt bās thāmālen.}
\end{quote}

\begin{quote}
The bus driver suddenly stopped the bus.
\end{quote}

Beside \bn{চালানো}, to drive: the \bn{চালক} is the one who \bn{চালায়}.

\begin{cousinweb}
\textbf{\bn{চালক}} (\emph{chālôk}) is a Sanskrit-style formation built from \emph{cāl-}, the causing stem of \emph{cal}, \textquotedblleft{}to move, to stir\textquotedblright{} (Monier-Williams), with the ending -\bn{ক} for the one who does a thing: one who makes something go. Monier-Williams records \emph{cālaka} only in an unrelated old sense, so the driver is a modern formation, the noun beside \bn{চালানো}.
\end{cousinweb}

\begin{grammarlens}[title={What you've built}]
One more word for getting around.
\end{grammarlens}

\subsection*{Your turn}
\begin{itemize}
\raggedright
  \item \emph{Say it:} \emph{chālôk}
  \item \emph{Say it:} \emph{chālôk}, once more
  \item \emph{Say it:} \emph{bāser chālôk hôṭhāt bās thāmālen.} — The bus driver suddenly stopped the bus.
  \item \emph{Recall:} say the Bengali for a van, then the Bengali for to race, then say \emph{chālôk} again
\end{itemize}

\begin{tcolorbox}[breakable,colback=teal!4,colframe=teal!35!black,title={Before you move on}]
What does \bn{চালক} mean? (\textquotedblleft{}A driver\textquotedblright{}.) Say it once more.
\end{tcolorbox}
\section[\texorpdfstring{\bn{কন্ডাক্টর} (kônḍākṭôr)}{কন্ডাক্টর (kônḍākṭôr)}]{\bn{কন্ডাক্টর} (kônḍākṭôr) — a bus conductor}

\label{lesson:BN-C389-kondaktor}



\begin{quote}
Before the new one: say the Bengali for to race, then the Bengali for a driver.
\end{quote}

\subsection*{You'll want to know: \bn{কন্ডাক্টর}}
\textbf{\bn{কন্ডাক্টর}} — \emph{kônḍākṭôr} — \textquotedblleft{}a bus conductor\textquotedblright{}.

\begin{quote}
\bn{কন্ডাক্টর} \bn{আমার} \bn{কাছে} \bn{ভাড়া} \bn{চাইলেন।}
\end{quote}

\begin{quote}
\emph{kônḍākṭôr āmār kāchhe bhāṛā chāilen.}
\end{quote}

\begin{quote}
The conductor asked me for the fare.
\end{quote}

\begin{cousinweb}
\textbf{\bn{কন্ডাক্টর}} (\emph{kônḍākṭôr}) is borrowed from English \emph{conductor}, with the English d and t said as the hard Bengali \bn{ড} and \bn{ট}.
\end{cousinweb}

\begin{grammarlens}[title={What you've built}]
One more word for getting around.
\end{grammarlens}

\subsection*{Your turn}
\begin{itemize}
\raggedright
  \item \emph{Say it:} \emph{kônḍākṭôr}
  \item \emph{Say it:} \emph{kônḍākṭôr}, once more
  \item \emph{Say it:} \emph{kônḍākṭôr āmār kāchhe bhāṛā chāilen.} — The conductor asked me for the fare.
  \item \emph{Recall:} say the Bengali for to race, then the Bengali for a driver, then say \emph{kônḍākṭôr} again
\end{itemize}

\begin{tcolorbox}[breakable,colback=teal!4,colframe=teal!35!black,title={Before you move on}]
What does \bn{কন্ডাক্টর} mean? (\textquotedblleft{}A bus conductor\textquotedblright{}.) Say it once more.
\end{tcolorbox}
\section[\texorpdfstring{\bn{হেলমেট} (helmeṭ)}{হেলমেট (helmeṭ)}]{\bn{হেলমেট} (helmeṭ) — a helmet}

\label{lesson:BN-C389-helmet}



\begin{quote}
Before the new one: say the Bengali for a driver, then the Bengali for a bus conductor.
\end{quote}

\subsection*{You'll want to know: \bn{হেলমেট}}
\textbf{\bn{হেলমেট}} — \emph{helmeṭ} — \textquotedblleft{}a helmet\textquotedblright{}.

\begin{quote}
\bn{মোটরসাইকেলে} \bn{হেলমেট} \bn{পরতে} \bn{হবে।}
\end{quote}

\begin{quote}
\emph{moṭôrsāikele helmeṭ pôrte hôbe.}
\end{quote}

\begin{quote}
On a motorbike you have to wear a helmet.
\end{quote}

\begin{cousinweb}
\textbf{\bn{হেলমেট}} (\emph{helmeṭ}) is borrowed from English \emph{helmet}, with the final t said as \bn{ট}.
\end{cousinweb}

\begin{grammarlens}[title={What you've built}]
One more word for getting around.
\end{grammarlens}

\subsection*{Your turn}
\begin{itemize}
\raggedright
  \item \emph{Say it:} \emph{helmeṭ}
  \item \emph{Say it:} \emph{helmeṭ}, once more
  \item \emph{Say it:} \emph{moṭôrsāikele helmeṭ pôrte hôbe.} — On a motorbike you have to wear a helmet.
  \item \emph{Recall:} say the Bengali for a driver, then the Bengali for a bus conductor, then say \emph{helmeṭ} again
\end{itemize}

\begin{tcolorbox}[breakable,colback=teal!4,colframe=teal!35!black,title={Before you move on}]
What does \bn{হেলমেট} mean? (\textquotedblleft{}A helmet\textquotedblright{}.) Say it once more.
\end{tcolorbox}
\section[\texorpdfstring{\bn{ইশারা} \bn{করা} (ishārā kôrā)}{ইশারা করা (ishārā kôrā)}]{\bn{ইশারা} \bn{করা} (ishārā kôrā) — to signal; to make a sign}

\label{lesson:BN-C389-ishara-kora}



\begin{quote}
Before the new one: say the Bengali for a bus conductor, then the Bengali for a helmet.
\end{quote}

\subsection*{You'll want to know: \bn{ইশারা} \bn{করা}}
\textbf{\bn{ইশারা} \bn{করা}} — \emph{ishārā kôrā} — \textquotedblleft{}to signal; to make a sign\textquotedblright{}.

\begin{quote}
\bn{কন্ডাক্টর} \bn{হাত} \bn{তুলে} \bn{চালককে} \bn{থামার} \bn{ইশারা} \bn{করলেন।}
\end{quote}

\begin{quote}
\emph{kônḍākṭôr hāt tule chālôkke thāmār ishārā kôrlen.}
\end{quote}

\begin{quote}
The conductor raised his hand and signalled the driver to stop.
\end{quote}

Past (he/she …): \bn{তিনি} \bn{চালককে} \bn{ইশারা} \bn{করলেন}, \emph{tini chālôkke ishārā kôrlen}.

The noun alone is \bn{ইশারা}, a sign or a hint: \bn{চোখের} \bn{ইশারা}, a sign with the eyes.

\begin{cousinweb}
\textbf{\bn{ইশারা} \bn{করা}} (\emph{ishārā kôrā}) joins \bn{ইশারা}, \textquotedblleft{}a sign, a gesture\textquotedblright{}, usually traced to Arabic \emph{ishāra}, \textquotedblleft{}a pointing, a signal\textquotedblright{}, through Persian, to \bn{করা}, to do.
\end{cousinweb}

\begin{grammarlens}[title={What you've built}]
One more word for getting around.
\end{grammarlens}

\subsection*{Your turn}
\begin{itemize}
\raggedright
  \item \emph{Say it:} \emph{ishārā kôrā}
  \item \emph{Say it:} \emph{ishārā kôrā}, once more
  \item \emph{Say it:} \emph{kônḍākṭôr hāt tule chālôkke thāmār ishārā kôrlen.} — The conductor raised his hand and signalled the driver to stop.
  \item \emph{Recall:} say the Bengali for a bus conductor, then the Bengali for a helmet, then say \emph{ishārā kôrā} again
\end{itemize}

\begin{tcolorbox}[breakable,colback=teal!4,colframe=teal!35!black,title={Before you move on}]
What does \bn{ইশারা} \bn{করা} mean? (\textquotedblleft{}To signal; to make a sign\textquotedblright{}.) Say it once more.
\end{tcolorbox}
\section[\texorpdfstring{\bn{হাঁকা} (hā̃kā)}{হাঁকা (hā̃kā)}]{\bn{হাঁকা} (hā̃kā) — to call out; to cry (stops or wares)}

\label{lesson:BN-C389-haka}



\begin{quote}
Before the new one: say the Bengali for a helmet, then the Bengali for to signal.
\end{quote}

\subsection*{You'll want to know: \bn{হাঁকা}}
\textbf{\bn{হাঁকা}} — \emph{hā̃kā} — \textquotedblleft{}to call out; to cry (stops or wares)\textquotedblright{}.

\begin{quote}
\bn{কন্ডাক্টর} \bn{পাড়ার} \bn{নাম} \bn{হাঁকলেন।}
\end{quote}

\begin{quote}
\emph{kônḍākṭôr pāṛār nām hā̃klen.}
\end{quote}

\begin{quote}
The conductor called out the name of the neighbourhood.
\end{quote}

Past (he/she …): \bn{তিনি} \bn{নাম} \bn{হাঁকলেন}, \emph{tini nām hā̃klen}.

Conductors \bn{হাঁকে} the stops from the door; hawkers \bn{হাঁকে} their goods in the lane.

\begin{cousinweb}
\textbf{\bn{হাঁকা}} (\emph{hā̃kā}) is inherited: CDIAL lists Bengali \emph{hā̃kā}, \textquotedblleft{}to call, to drive off\textquotedblright{}, under Sanskrit \emph{hakkayati}, \textquotedblleft{}calls to\textquotedblright{} (CDIAL 13939).
\end{cousinweb}

\begin{grammarlens}[title={What you've built}]
One more word for getting around.
\end{grammarlens}

\subsection*{Your turn}
\begin{itemize}
\raggedright
  \item \emph{Say it:} \emph{hā̃kā}
  \item \emph{Say it:} \emph{hā̃kā}, once more
  \item \emph{Say it:} \emph{kônḍākṭôr pāṛār nām hā̃klen.} — The conductor called out the name of the neighbourhood.
  \item \emph{Recall:} say the Bengali for a helmet, then the Bengali for to signal, then say \emph{hā̃kā} again
\end{itemize}

\begin{tcolorbox}[breakable,colback=teal!4,colframe=teal!35!black,title={Before you move on}]
What does \bn{হাঁকা} mean? (\textquotedblleft{}To call out; to cry (stops or wares)\textquotedblright{}.) Say it once more.
\end{tcolorbox}
